\documentclass[a4paper,11pt]{article}
\usepackage[utf8]{inputenc}
\usepackage[spanish]{babel}
\usepackage{geometry}
\usepackage{enumitem}
\usepackage{sectsty}
\usepackage{titlesec}
\usepackage{tikz}
\usepackage{array}

\geometry{
    a4paper,
    left=20mm,
    top=18mm,
    right=20mm,
    bottom=20mm
}

% Sin números de página
\pagestyle{empty}

% Estilo de secciones con números en círculos
\newcommand{\circled}[1]{%
    \tikz[baseline=(char.base)]{%
        \node[shape=circle,draw,inner sep=0.5pt,fill=black,text=white,font=\bfseries\footnotesize] (char) {#1};%
    }%
}

% Formato de secciones
\titleformat{\section}
{\normalfont\Large\bfseries}
{\circled{\thesection}}{0.5em}{}

\titlespacing*{\section}{0pt}{12pt}{8pt}

\titleformat{\subsection}
{\normalfont\normalsize\bfseries}
{}{0em}{}

\titlespacing*{\subsection}{8pt}{4pt}{4pt}

\begin{document}

% Encabezado del curso
\begin{center}
\small{Universidad de Granada}

\vspace{0.3cm}

\textbf{\large{Producción e Interacción Oral - Nivel 7}}

\vspace{0.4cm}

\small
\begin{tabular}{@{}ll@{}}
\textbf{Grupos:} & A (Aula 24) y B (Aula 08) \\
\textbf{Centro:} & Centro de Lenguas Modernas / Edificio A \\
\textbf{Profesor:} & Javier Benítez Lainez \\
\textbf{Tutorías:} & Martes 9:00-10:00 y 14:30-15:30 \\
& Sala de profesores / Email: benitezl@go.ugr.es
\end{tabular}
\end{center}

\vspace{0.5cm}
\section*{Ejercicios de Narración Temporal}

\vspace{0.4cm}

\section{Práctica de Pretéritos e Indicadores Temporales C1}

\vspace{0.2cm}

---

\vspace{0.2cm}

\section{INTRODUCCIÓN}

\vspace{0.2cm}

Este cuaderno contiene ejercicios prácticos para dominar la narración temporal en español nivel C1.

\vspace{0.2cm}

---

\vspace{0.2cm}

\section{EJERCICIO 1: Selección de Tiempo}

\vspace{0.2cm}

Elige pretérito indefinido o imperfecto:

\vspace{0.2cm}

1. "Ayer ____ (hacía/hizo) mucho calor."

\textbf{Respuesta:} hizo

\vspace{0.2cm}

2. "____ (Era/Fue) una noche oscura y ____ (llovía/llovió)."

\textbf{Respuesta:} Era, llovía

\vspace{0.2cm}

3. "Cuando ____ (era/fui) niño, ____ (iba/fui) al colegio todos los días."

\textbf{Respuesta:} era, iba

\vspace{0.2cm}

4. "Mientras ____ (cocinaba/cociné), ____ (escuchaba/escuché) música."

\textbf{Respuesta:} cocinaba, escuchaba

\vspace{0.2cm}

5. "____ (Tenía/Tuve) un gato que se ____ (llamaba/llamó) Michi."

\textbf{Respuesta:} Tenía, llamaba

\vspace{0.2cm}

---

\vspace{0.2cm}

\section{EJERCICIO 2: Completa la Historia}

\vspace{0.2cm}

Completa con el tiempo apropiado (indefinido/imperfecto/pluscuamperfecto):

\vspace{0.2cm}

"Cuando ____ (ser) joven, ____ (vivir) en Madrid. Todos los días ____ (despertarse) a las 7 y ____ (tomar) desayuno. Luego ____ (ir) al trabajo. Un día, mientras ____ (trabajar), ____ (recibir) una noticia importante. No ____ (esperar) que ____ (ser) tan relevante. Ya ____ (hablar) con mi jefe sobre el tema anteriormente, pero esa vez ____ (ser) diferente. Inmediatamente ____ (tomar) una decisión que ____ (cambiar) mi vida."

\vspace{0.2cm}

\textbf{Respuesta:} era, vivía, me despertaba, tomaba, iba, trabajaba, recibí, esperaba, era, había hablado, fue, tomé, cambiaría/cambió

\vspace{0.2cm}

---

\vspace{0.2cm}

\section{EJERCICIO 3: Uso de Conectores}

\vspace{0.2cm}

Ordena los eventos usando conectores:

\vspace{0.2cm}

Eventos:

\begin{itemize}

  \item Terminé el proyecto

  \item Empezaron las dificultades

  \item Finalmente lo logré

  \item Al principio fue fácil

  \item Luego me di cuenta de que era difícil

\end{itemize}

\vspace{0.2cm}

\textbf{Respuesta:} "Al principio fue fácil. Luego me di cuenta de que era difícil y empezaron las dificultades. Sin embargo, finalmente lo logré después de terminar el proyecto."

\vspace{0.2cm}

---

\vspace{0.2cm}

\section{EJERCICIO 4: Narración Personal}

\vspace{0.2cm}

Escribe una breve narración (5-8 oraciones) sobre un evento memorable de tu pasado usando:

\begin{itemize}

  \item Al menos 3 conectores temporales diferentes

  \item Al menos 3 tiempos pasados diferentes

  \item Al menos 2 indicadores de frecuencia o duración

\end{itemize}

\vspace{0.2cm}

\textbf{Plantilla:}

"En aquella época... ____ (indicador de frecuencia/época), ____ (imperfecto)______. Un día ____ (indefinido)______. Mientras ____ (imperfecto)______, ____ (pluscuamperfecto)______. Finalmente ____ (indefinido)______."

\vspace{0.2cm}

---

\vspace{0.2cm}

\section{EJERCICIO 5: Corrección de Errores}

\vspace{0.2cm}

Corrige los errores en estas oraciones:

\vspace{0.2cm}

1. "Ayer he hablado con María."

\textbf{Corrección:} Ayer hablé con María.

\vspace{0.2cm}

2. "Cuando era niño, fui al colegio todos los días."

\textbf{Corrección:} Cuando era niño, iba al colegio todos los días.

\vspace{0.2cm}

3. "Esta semana comí en ese restaurante."

\textbf{Corrección:} Esta semana he comido en ese restaurante.

\vspace{0.2cm}

4. "Llegué y ya salió."

\textbf{Corrección:} Llegué y ya había salido.

\vspace{0.2cm}

5. "Mientras cocinaba, escuchó las noticias."

\textbf{Corrección:} Mientras cocinaba, escuchaba las noticias.

\vspace{0.2cm}

---

\vspace{0.2cm}

\section{EJERCICIO 6: Preguntas de Narración}

\vspace{0.2cm}

Responde estas preguntas narrando en pasado:

\vspace{0.2cm}

1. ¿Qué hiciste ayer?

\textbf{Respuesta modelo:} "Ayer me levanté a las 8, desayuné y fui a trabajar."

\vspace{0.2cm}

2. ¿Cómo era tu rutina cuando eras niño?

\textbf{Respuesta modelo:} "Cuando era niño, me despertaba temprano, iba al colegio y por la tarde jugaba."

\vspace{0.2cm}

3. ¿Qué has hecho esta semana?

\textbf{Respuesta modelo:} "Esta semana he trabajado mucho y he ido al gimnasio varias veces."

\vspace{0.2cm}

---

\vspace{0.2cm}

\section{EJERCICIO 7: Role-Play Narrativo}

\vspace{0.2cm}

En parejas, narrar un evento desde perspectivas diferentes. Usa conectores y tiempos apropiados.

\vspace{0.2cm}

\textbf{Escenario:} Un viaje que salió mal

\vspace{0.2cm}

\textbf{Persona A (el narrador optimista):} Narrar destacando los aspectos positivos

\vspace{0.2cm}

\textbf{Persona B (el narrador pesimista):} Narrar destacando los problemas

\vspace{0.2cm}

\textbf{Conectores útiles:} al principio, luego, sin embargo, afortunadamente, desafortunadamente, por último, etc.

\vspace{0.2cm}

---

\vspace{0.2cm}

\section{EJERCICIO 8: Transformación de Tiempos}

\vspace{0.2cm}

Transforma la historia de presente a pasado:

\vspace{0.2cm}

"Vivo en Madrid. Todos los días tomo el metro a las 8. Mientras voy en el metro, leo el periódico. Un día, conozco a alguien interesante. Nos hacemos amigos."

\vspace{0.2cm}

\textbf{Transformación:} "Vivía en Madrid. Todos los días tomaba el metro a las 8. Mientras iba en el metro, leía el periódico. Un día, conocí a alguien interesante. Nos hicimos amigos."

\vspace{0.2cm}

---

\vspace{0.2cm}

\textbf{Sesión 11: Ejercicios de Narración Temporal}

\textbf{Nivel C1 - Español Oral Avanzado}

\vspace{0.2cm}

\end{document}
